% !TeX root=tukedip.tex
\section{Záver (zhodnotenie riešenia)}
\label{sec:conclusion}

Práca sa zaoberala návrhom, implementáciou a~experimentálnym hodnotením bezpečnostných mechanizmov pre systém kontroly fyzického prístupu na báze RFID/NFC. Výsledkom je funkčný prototyp pozostávajúci z~hardvérovej čítačky (ESP32 + RC522) a~dvoch nezávislých serverových služieb v~jazyku C++20 --- frontend služba (spracovanie požiadaviek, HMAC autentifikácia, konštrukcia AEAD frame) a~DecisionService (dešifrovanie, rozhodovacia logika, anomálna detekcia) --- s~lokálnym úložiskom SQLite. Systém implementuje viacvrstvovú ochranu: HMAC-SHA-256 autentifikáciu hardvérových požiadaviek, AEAD frame (ChaCha20-Poly1305 / XChaCha20-Poly1305) s~AAD binding, HKDF per-reader deriváciu kľúčov, anti-replay sliding window, anomálny detektor s~karanténou a~tamper-evident audit log. Architektúra s~oddeleným DecisionService umožňuje nezávislé testovanie bezpečnostných vlastností AEAD vrstvy; scenár S7 (simulovaný nonce tamper) modeluje efekt MITM útoku na frame pred vstupom do rozhodovacej vrstvy.

Okrem samotného prototypu práca prináša reprodukovateľný experimentálny rámec: automatizovaný harness izoluje útočný vstup, cieľové primitívy a~metriky do samostatných komponentov, Python pipeline agreguje výstupy do CSV a~generuje grafy aj LaTeX tabuľky použité v~texte. Porovnanie profilov prebieha nad spoločnou produkčnou kódovou základňou so zdieľaným AEAD enginom, čím je zabezpečená neutralita vyhodnotenia --- žiadny profil nemá implementačnú výhodu nad ostatnými.

Prototyp je kontajnerizovateľný pomocou priloženého \texttt{Dockerfile}; modulárna štruktúra serverovej aplikácie umožňuje výmenu jednotlivých vrstiev bez dopadu na ostatné (napr.\ nahradenie SQLite plnohodnotnou SQL databázou alebo pridanie ďalších AEAD algoritmov). Použitie výlučne permisívnych open-source závislostí odstraňuje licenčné prekážky pre ďalšie rozšírenia.

Experimentálna analýza šiestich profilov (tabuľka~\ref{tab:findings}) (A1/A2 $\times$ R0/R1/R2) v~siedmich scenároch potvrdila hlavné zistenia:

\begin{table}[H]
\renewcommand{\arraystretch}{1.3}
\centering
\caption{Súhrn hlavných experimentálnych zistení}
\label{tab:findings}
\small
\begin{tabular}{c p{10.5cm}}
\toprule
\textbf{Tézis} & \textbf{Zistenie} \\
\midrule
T1 & R2 ako jediný vynucuje nonce politiku (S3a: R0/R1 100\% úspešnosť vs.\ R2 0\%). HKDF per-reader derivácia zabraňuje cross-reader útokom (S3b: 0\% naprieč všetkými profilmi). \\
T2 & \texttt{ProtocolAnomalyDetector} (profily R2) dosiahol 0\% úspešnosť naprieč všetkými testovanými scenármi (S1--S7). V~S3a a~S4 bol jediným mechanizmom, ktorý scenár zastavil; v~S7 detekoval nonce tamper ešte pred pokusom o~dešifrovanie. \\
T3 & Replay window a~anomálna detekcia sú komplementárne. S4: 11\% pre R0/R1 $\rightarrow$ 0\% pre R2. \\
T4 & AAD binding poskytuje ochranu nezávisle od režimu nonce (S2: 0\% naprieč všetkými profilmi). S7 potvrdzuje, že modifikácia nonce v~zostavenom frame vedie k~zlyhaniu AEAD tag verifikácie aj bez detektora. \\
T5 & Karanténa mení reakciu z~per-frame odmietnutia na izoláciu zariadenia (defense-in-depth). \\
T6 & Overhead detektora (${\sim}15$--$19\%$, ${\sim}25$--$35\,\mu$s na frame v~závislosti od AEAD algoritmu) je rádovo nepodstatný pre reálnu prevádzku ($<$10 udalostí/s), hoci nie je zanedbateľný. \\
\bottomrule
\end{tabular}
\end{table}

\noindent\textbf{Obmedzenia a~budúca práca}

Práca vedome zjednodušuje niektoré aspekty produkčného nasadenia: karta slúži ako identifikátor, TLS nie je povinné a~správa tajomstiev je pragmatická (hex súbor, nie HSM). Tieto obmedzenia vymedzujú rámec pre interpretáciu experimentov.

Smery budúcej práce: mTLS\cite{RFC8446} pre transport, migrácia na kryptografické smart karty (napr.\ MIFARE DESFire\cite{GarciaMIFARE2008}) a~integrácia s~centralizovanou správou tajomstiev (HSM/vault)\cite{NIST80057}.
